\begingroup
\makeatletter\tagsleft@true\let\veqno\@@leqno\makeatother
\ArithmeticHead{345}{R}{{\scriptsize ZUR THEORIE DER REIHE $F(\alpha,\beta,\gamma,x)$.\quad III.}}
\NumberedItem{5}
\noindent {\small [S. 42]}\par
Mit Hülfe des vorhergehenden hat man nun auch sofort das Nöthige zur Bestimmung der Coefficienten der Reihe, in welche
\[\frac{1}{\surd(aa+a'a'-2aa'\cos\varphi)}\]
\noindent entwickelt wird. Setzt man dieselbe
\[T=A^0+2A'\cos\varphi+2A''\cos2\varphi+2A'''\cos3\varphi\ \text{u.s.w.},\]
\noindent so wird, $f$ statt $\frac{a}{a'}$ geschrieben und $r$[*)] statt $\cos\varphi+i\sin\varphi$,
\[\frac{1}{a'}(1-fr)^{-\frac12}(1-fr^{-1})^{-\frac12}=T,\]
\noindent also
\[
\begin{aligned}
a'A^0={}&1+\frac{1{.}1}{2{.}2}ff+\frac{1{.}1{.}3{.}3}{2{.}2{.}4{.}4}f^4
 +\frac{1{.}1{.}3{.}3{.}5{.}5}{2{.}2{.}4{.}4{.}6{.}6}f^6\,{.}\,{.}\,{.}\\
={}&F\left(\frac12,\frac12,1,f^2\right),
\end{aligned}
\]
\[
\begin{aligned}
a'A'={}&\frac12 f\Bigl(1+\frac{1{.}3}{2{.}4}ff+\frac{1{.}3{.}3{.}5}{2{.}4{.}4{.}6}f^4\\
&\qquad+\frac{1{.}3{.}3{.}5{.}5{.}7}{2{.}4{.}4{.}6{.}6{.}8}f^6\,{.}\,{.}\,{.}\Bigr)\\
={}&\frac12 f{.}F\left(\frac12,\frac32,2,f^2\right),
\end{aligned}
\]
\[
\begin{aligned}
a'A''={}&\frac12{.}\frac34 ff\Bigl(1+\frac{1{.}5}{2{.}6}ff+\frac{1{.}5{.}3{.}7}{2{.}6{.}4{.}8}f^4\\
&\qquad+\frac{1{.}5{.}3{.}7{.}5{.}9}{2{.}6{.}4{.}8{.}6{.}10}f^6\,{.}\,{.}\,{.}\Bigr)\\
={}&\frac{1{.}3}{2{.}4}ff{.}F\left(\frac12,\frac52,3,ff\right),
\end{aligned}
\]
\[
\begin{aligned}
a'A'''={}&\frac{1{.}3{.}5}{2{.}4{.}6}f^3\Bigl(1+\frac{1{.}7}{2{.}8}ff+\frac{1{.}7{.}3{.}9}{2{.}8{.}4{.}10}f^4\\
&\qquad+\frac{1{.}7{.}3{.}9{.}5{.}11}{2{.}8{.}4{.}10{.}6{.}12}f^6\,{.}\,{.}\,{.}\Bigr)\\
={}&\frac{1{.}3{.}5}{2{.}4{.}6}f^3{.}F\left(\frac12,\frac72,4,ff\right)
\end{aligned}
\]
\centerline{u.s.w.}
\par\medskip\noindent\rule{.25\linewidth}{.4pt}\par\smallskip
{\small\noindent [*) Statt $r$ steht in der Handschrift $\rho$.]}
\par\medskip\noindent X 1.\hfill 44
\clearpage\SourceLeaf{353}\ArithmeticHead{346}{L}{ANALYSIS.\quad NACHLASS.}
Aus der Verbindung der Gleichungen (15) und (19) folgt
\[
\begin{aligned}
&\beta(\alpha-\gamma)xF(\alpha,\beta+1,\gamma+1)\\
&\quad-\gamma(\gamma-1)F(\alpha,\beta-1,\gamma-1)\\
&\quad+\gamma(\gamma-1+(\beta-\alpha)x)F(\alpha,\beta,\gamma)=0.
\end{aligned}
\]
\noindent Also für unsern Fall, wo $\alpha=\frac12$, $\beta=\gamma-\frac12$,
\[
\begin{aligned}
&-\left(\gamma-\frac12\right)^2xF\left(\frac12,\gamma+\frac12,\gamma+1\right)\\
&\quad-\gamma(\gamma-1)F\left(\frac12,\gamma-\frac32,\gamma-1\right)\\
&\quad+\gamma(\gamma-1+(\gamma-1)x)F\left(\frac12,\gamma-\frac12,\gamma\right)=0.
\end{aligned}
\]
\noindent Es ist folglich, weil
\[ffA^{(\gamma-2)},\qquad fA^{(\gamma-1)};\qquad A^{(\gamma)}\]
\noindent proportional sind
\[
\begin{gathered}
\gamma(\gamma-1)F\left(\frac12,\gamma-\frac32,\gamma-1\right),\\
\left(\gamma-\frac32\right)\gamma F\left(\frac12,\gamma-\frac12,\gamma\right),\\
\left(\gamma-\frac32\right)\left(\gamma-\frac12\right)F\left(\frac12,\gamma+\frac12,\gamma+1\right),
\end{gathered}
\]
\noindent allgemein[*)]
\begin{equation}
\begin{aligned}
&\left(\gamma-\frac32\right)A^{(\gamma-2)}-(\gamma-1)\left(f+\frac1f\right)A^{(\gamma-1)}\\
&\quad+\left(\gamma-\frac12\right)A^{(\gamma)}=0,
\end{aligned}\tag*{(43)}
\end{equation}
\noindent also
\[
\begin{aligned}
A^{(0)}-2\left(f+\frac1f\right)A'+3A''&=0,\\
3A'-4\left(f+\frac1f\right)A''+5A'''&=0,\\
5A''-6\left(f+\frac1f\right)A'''+7A^{\mathrm{IV}}&=0\ \mathrm{etc.}
\end{aligned}
\]
\noindent Man setze
\begin{equation}
\begin{gathered}
\frac21\cdot\frac{1+ff}{f}\cdot\frac{A'}{A^0}=B^0,\qquad
\frac43\cdot\frac{1+ff}{f}\cdot\frac{A''}{A'}=B',\\
\frac65\cdot\frac{1+ff}{f}\cdot\frac{A'''}{A''}=B''\ \mathrm{etc.},
\end{gathered}\tag*{(44)}
\end{equation}
\noindent so wird
\begin{equation}
\begin{gathered}
B^0=\frac{1}{1-\frac98\left(\frac{f}{1+ff}\right)^2B'},\qquad
B'=\frac{1}{1-\frac{25}{24}\left(\frac{f}{1+ff}\right)^2B''},\\[1.2ex]
B''=\frac{1}{1-\frac{49}{48}\left(\frac{f}{1+ff}\right)^2B'''}\\
\mathrm{etc.}
\end{gathered}\tag*{(45)}
\end{equation}
\noindent {\small [S. 43]}\par
Nun ist aber auch
\par\medskip\noindent\rule{.25\linewidth}{.4pt}\par\smallskip
{\small\noindent [*) In der Handschrift steht hier »und«.]}\par
\clearpage\SourceLeaf{354}\ArithmeticHead{347}{R}{{\scriptsize ZUR THEORIE DER REIHE $F(\alpha,\beta,\gamma,x)$.\quad III.}}
\[
\begin{gathered}
B^0=\frac{(1+ff)F\left(\frac12,\frac32,2,ff\right)}{F\left(\frac12,\frac12,1,ff\right)},\\[1ex]
B'=\frac{(1+ff)F\left(\frac12,\frac52,3,ff\right)}{F\left(\frac12,\frac32,2,ff\right)}\ \mathrm{etc.}
\end{gathered}
\]
\noindent Man hat also auch
\begin{equation}
B^{(n-2)}=\cfrac[l]{1+ff}{1-\cfrac[l]{\dfrac{1{.}1}{(2n-2)2n}ff}
 {1-\cfrac[l]{\dfrac{(2n-1)^2}{2n(2n+2)}ff}
 {1-\cfrac[l]{\dfrac{9}{(2n+2)(2n+4)}ff}
 {1-\cfrac[l]{\dfrac{(2n+1)^2}{(2n+4)(2n+6)}ff}
 {1+\mathrm{etc.}}}}}}
\tag*{(46)}
\end{equation}
\clearpage
\endgroup
